% !PRINTED 33
\weberpage{33}

Für die Theilbarkeit von Functionen gelten dieselben Gesetze, wie für die Theilbarkeit der Zahlen, insbesondere die folgenden:

\gesperrt{1. Wenn die Function $F(x)$ durch die Function $f(x)$, $f(x)$ durch eine dritte Function $\varphi(x)$ theilbar ist, so ist auch $F(x)$ durch $\varphi(x)$ theilbar.}

Denn ist
\[
 F=Qf,
 \qquad
 f=q\varphi,
\]
worin $Q$ und $q$ ganze rationale Functionen sind, so ist
\[
 F=Qq\varphi,
\]
und da $Qq$ eine ganze rationale Function ist, $F$ durch $\varphi$ theilbar.

\gesperrt{2. Ist $F(x)$ durch $f(x)$ theilbar und $Q$ eine beliebige ganze rationale Function, so ist auch $QF(x)$ durch $f(x)$ theilbar.}

\gesperrt{3. Ist $F(x)$ und $f(x)$ durch $\varphi(x)$ theilbar, so ist auch $F(x)\pm f(x)$ durch $\varphi(x)$ theilbar, oder allgemeiner:}

\gesperrt{4. Sind $F_1,F_2\ldots$ durch $f(x)$ theilbar und $Q_1,Q_2\ldots$ beliebige ganze rationale Functionen, so ist auch}
\[
 Q_1F_1+Q_2F_2+\cdots
\]
\gesperrt{durch $f(x)$ theilbar.}

Der letzte Satz umfasst die beiden vorhergehenden und wird einfach so bewiesen.

Sind $F_1,F_2\ldots$ durch $f$ theilbar, so kann man die ganzen rationalen Functionen $\Phi_1,\Phi_2\ldots$ so bestimmen, dass
\[
 F_1=\Phi_1f,\qquad F_2=\Phi_2f,\ldots
\]
und folglich
\[
 Q_1F_1+Q_2F_2+\cdots=(Q_1\Phi_1+Q_2\Phi_2+\cdots)f.
\]
Da nun $Q_1\Phi_1+Q_2\Phi_2+\cdots$ eine ganze rationale Function ist, so ist der Satz bewiesen.

\gesperrt{5. Jede Function ist durch sich selbst theilbar.}

In Bezug auf die Theilbarkeit oder Untheilbarkeit von Functionen wird nichts geändert, wenn die Functionen mit beliebigen, von $x$ unabhängigen Factoren multiplicirt werden.

Eine von $x$ unabhängige, von Null verschiedene Grösse kann als Function nullten Grades aufgefasst werden. Nennen wir eine solche Grösse eine Constante, so können wir sagen:

\gesperrt{6. Jede Function ist durch jede Constante theilbar.}
